% Auxiliary common-alphabet forms of SCP10 eq:2d:closure-inv and eq:2d:close-in-in.
\subsection{Physical inverses and matched projector contractions}
\label{sec:peps_four_cut_common_projectors}

\begin{definition}[Products of a family of physical maps]%
    \label{def:peps_four_cut_common_physical_product}
    \lean{TNLean.PEPS.physicalProductFamilyMatrix,
        TNLean.PEPS.physicalProductFamilyMap,
        TNLean.PEPS.physicalProductFamilyMatrix_const,
        TNLean.PEPS.physicalProductFamilyMap_const,
        TNLean.PEPS.physicalProductFamilyMap_apply,
        TNLean.PEPS.torusSitewisePhysicalMap}
    \leanok
    \uses{def:peps_physical_product_map}
    For a finite site set $I$ and matrices $F_v:\C^P\to\C^Q$, put
    \begin{align}
        (\mathcal T_F)_{\tau,\sigma} & =\prod_{v\in I}F_v(\tau_v,\sigma_v),
        \nonumber                                                                                           \\
        (\mathcal T_F\psi)(\tau)
                                     & =\sum_{\sigma:I\to P}\prod_{v\in I}F_v(\tau_v,\sigma_v)\psi(\sigma).
        \label{eq:peps_common_physical_product}
    \end{align}
    The map formula assumes $P$ finite; the matrix formula does not.
    A constant family gives $F^{\otimes I}$. Taking $I=X$ gives the
    site-dependent physical map on the torus.
\end{definition}

\begin{theorem}[Composition and coordinate action of product maps]%
    \label{thm:peps_four_cut_common_product_algebra}
    \lean{TNLean.PEPS.physicalProductFamilyMatrix_mul,
        TNLean.PEPS.physicalProductFamilyMap_comp,
        TNLean.PEPS.physicalProductFamilyMatrix_one,
        TNLean.PEPS.physicalProductFamilyMap_one,
        TNLean.PEPS.physicalProductFamilyMap_oneCoordinate_apply,
        TNLean.PEPS.physicalProductFamilyMap_fixed_of_slices_fixed}
    \leanok
    \uses{def:peps_four_cut_common_physical_product}
    For compatible matrices $F_v,L_v$ with finite intermediate alphabets,
    $\mathcal T_F\mathcal T_L=\mathcal T_{(F_vL_v)_v}$ as matrices;
    the same equality holds as maps when the input alphabet is finite.
    A product of identity matrices is the identity matrix and, for finite
    alphabets, the identity map. If $P$ is finite, $B$ is a matrix on
    $\C^P$, and $F_v=B$ at a single site $v$ and $F_w=\Id$ elsewhere, then
    \begin{align}
        (\mathcal T_F\psi)(\tau)
         & =\sum_{s\in P}B(\tau_v,s)\psi(\tau[v\mapsto s]).
        \label{eq:peps_common_one_coordinate}
    \end{align}
    In particular, if matrices $B_v$ fix every slice
    $s\mapsto\psi(\tau[v\mapsto s])$ at their respective sites, then
    $\mathcal T_B\psi=\psi$.
\end{theorem}
\begin{proof}\leanok
    \uses{thm:peps_dependent_product_calculus}
    Specialize the product identities to constant input and output
    alphabets. Finite distributivity factors the intermediate-coordinate
    sum into $F_vL_v$ at every site. Identity entries force agreement at
    all unchanged coordinates, giving (\ref{eq:peps_common_one_coordinate}).
    Apply the single-coordinate maps successively over the finite site set;
    every step fixes $\psi$, and their product is $\mathcal T_B$.
\end{proof}

\begin{theorem}[Slice characterization of a product range]%
    \label{thm:peps_four_cut_common_product_range}
    \lean{TNLean.PEPS.physicalProductFamilyMap_slice_mem_range,
        TNLean.PEPS.mem_range_physicalProductFamilyMap_iff}
    \leanok
    \uses{def:peps_four_cut_common_physical_product}
    Let $I,P$ be finite. Every one-coordinate slice of $\mathcal T_F\chi$
    lies in the corresponding range $\operatorname{ran}F_v$.
    If also $Q$ is finite, then
    \begin{align}
        \psi\in\operatorname{ran}\mathcal T_F
        \quad\Longleftrightarrow\quad
        \forall v\in I,\ \forall\tau:I\to Q,\quad
        \psi(\tau[v\mapsto\cdot])\in\operatorname{ran}F_v.
        \label{eq:peps_common_product_range}
    \end{align}
    Empty alphabets and an empty site set are allowed, as are zero maps
    and maps of different ranks.
\end{theorem}
\begin{proof}\leanok
    \uses{thm:peps_dependent_product_range}
    Apply Theorem~\ref{thm:peps_dependent_product_range} with constant
    alphabets. Explicitly, a slice of a product image is a linear combination
    of the columns of $F_v$. Conversely, choose $L_v$ with $F_vL_v$
    equal to the identity on $\operatorname{ran}F_v$.
    The slice conditions imply $\mathcal T_{FL}\psi=\psi$, so
    $\psi=\mathcal T_F(\mathcal T_L\psi)$.
\end{proof}

\begin{theorem}[Sitewise physical maps preserve the cut boundary]%
    \label{thm:peps_four_cut_common_physical_contraction}
    \lean{TNLean.PEPS.torusSitewisePhysicalMap_torusCutMap,
        TNLean.PEPS.torusSitewisePhysicalMap_torusBondNetwork,
        TNLean.PEPS.torusSitewisePhysicalMap_sitewiseTorusGClosure,
        TNLean.PEPS.torusSitewisePhysicalMap_matchedTorusGClosure,
        TNLean.PEPS.map_matchedCommutingClosureSpan}
    \leanok
    \uses{def:peps_four_cut_common_physical_product,
        def:peps_four_cut_common_boundary,def:peps_four_cut_common_closure}
    Let the input physical alphabet be finite, and set
    \begin{align}
        (FA)_v(t,r,b,l;q)=\sum_sF_v(q,s)A_v(t,r,b,l;s).
        \label{eq:peps_common_physical_site}
    \end{align}
    For every boundary $M$ and every pair of bond-operator families,
    \begin{align}
        \mathcal T_F\mathcal C^A_{c,r}M
                                      & =\mathcal C^{FA}_{c,r}M,
        \label{eq:peps_common_physical_cut}                                            \\
        \mathcal T_F\mathcal N_A(O^{\mathrm h},O^{\mathrm v})
                                      & =\mathcal N_{FA}(O^{\mathrm h},O^{\mathrm v}),
        \nonumber                                                                      \\
        \mathcal T_F\mathcal M_A(g,h) & =\mathcal M_{FA}(g,h),
        \nonumber                                                                      \\
        \mathcal T_F(\mathcal Z_A)    & =\mathcal Z_{FA}.
        \nonumber
    \end{align}
    The closure identities include uniform bond representations.
\end{theorem}
\begin{proof}\leanok
    Expand the physical and virtual sums and exchange their order. The
    physical sum factors into $\sum_sF_v(q,s)A_v(\cdot;s)$ at each site,
    leaving every boundary coefficient and bond operator untouched.
    Substitute the closure operators for the third equality, and take the
    image of the span of commuting closures for the fourth.
\end{proof}

\begin{theorem}[One local inverse for all four cuts]%
    \label{thm:peps_four_cut_common_local_inverse}
    \lean{TNLean.PEPS.physicalMapSite_representationAveragingSite,
        TNLean.PEPS.exists_sitewisePhysicalMap_to_representationAveragingSite,
        TNLean.PEPS.exists_torusCutMap_localInverse,
        TNLean.PEPS.torusCutMap_recover_representationAveragingSite,
        TNLean.PEPS.fourTorusCutSpace_eq_map_representationAveragingSite,
        TNLean.PEPS.map_matchedCommutingClosureSpan_recover}
    \leanok
    \uses{def:peps_four_cut_common_matched_action,
        def:peps_four_cut_common_physical_product,def:peps_four_cut_common_boundary,
        def:peps_four_cut_common_closure,def:peps_g_injective}
    Let $G$ be finite and let $\rho_v$ be arbitrary representations on
    $\C^{V^4}$. Put $R_v=\mathcal P(A_v)$.
    If $R_v\rho_v(k)=R_v$ for all $v,k$, then $R_v\Pi_{\rho_v}=R_v$,
    so $R_vA^0_{\rho_v}=A_v$ as tensors and
    \begin{align}
        \mathcal T_R\mathcal C^{A^0}_{c,r}M
         & =\mathcal C^A_{c,r}M.
        \label{eq:peps_common_cut_recovery}
    \end{align}
    If every $A_v$ is $G$-injective and the physical alphabet is finite,
    there exist matrices $L_v$ with $L_vR_v=\Pi_{\rho_v}$, hence
    $L_vA_v=A^0_{\rho_v}$ and, simultaneously for all $c,r,M$,
    \begin{align}
        \mathcal T_L\mathcal C^A_{c,r}M
         & =\mathcal C^{A^0}_{c,r}M.
        \label{eq:peps_common_cut_inverse}
    \end{align}
    On the two-by-two torus this gives
    $\mathcal S^A=\mathcal T_R(\mathcal S^{A^0})$.
    For the matched actions, invariance alone also gives
    $\mathcal T_R(\mathcal Z_{A^0})=\mathcal Z_A$.
    These are common-alphabet local-inverse identities supporting
    \cite[Theorem~5.5]{Schuch2010PEPS}.
\end{theorem}
\begin{proof}\leanok
    \uses{thm:peps_g_injective_left_inverse,
        thm:peps_site_coefficients_composition,
        thm:peps_four_cut_common_physical_contraction}
    Average $R_v\rho_v(k)=R_v$ to obtain $R_v\Pi_{\rho_v}=R_v$.
    Theorem~\ref{thm:peps_g_injective_left_inverse} supplies $L_v$
    independently at each site. Apply
    (\ref{eq:peps_common_physical_cut}) to these two local identities.
    If $\psi\in\mathcal S^A$, the same vector $\mathcal T_L\psi$
    lies in every canonical cut range by
    (\ref{eq:peps_common_cut_inverse}); recovering it using one cut gives
    $\mathcal T_R\mathcal T_L\psi=\psi$.
    Conversely, (\ref{eq:peps_common_cut_recovery}) maps each canonical
    range into the corresponding physical range. Finally, physical
    transport of the closure span and $R_vA^0_v=A_v$ give the asserted
    closure recovery.
\end{proof}

\begin{theorem}[Multilinear expansion of canonical cut contractions]%
    \label{thm:peps_four_cut_common_cut_expansion}
    \lean{TNLean.PEPS.torusCutCoeff_sum,
        TNLean.PEPS.torusCutCoeff_mul,
        TNLean.PEPS.representationAveragingSite_apply,
        TNLean.PEPS.torusCutCoeff_representationAveragingSite,
        TNLean.PEPS.representationAveragingSite_torusMatchedLegRep_apply}
    \leanok
    \uses{def:peps_four_cut_common_boundary,def:peps_four_cut_common_matched_action}
    For a finite family $A^i_v$ and scalars $z_v$, the cut contraction obeys
    \begin{align}
        \mathcal C^{(\sum_i A^i_v)_v}_{c,r}M
         & =\sum_{q:X\to I}\mathcal C^{(A^{q_v}_v)_v}_{c,r}M,
        \nonumber                                             \\
        \mathcal C^{(z_vA_v)_v}_{c,r}M
         & =\left(\prod_vz_v\right)\mathcal C^A_{c,r}M.
        \nonumber
    \end{align}
    For finite $G$, the canonical tensor has coefficients
    $A^0_{\rho_v}(c;s)=|G|^{-1}\sum_{k\in G}\rho_v(k)_{s,c}$.
    Consequently, if $B^q_v(c;s)=\rho_v(q_v)_{s,c}$, then
    \begin{align}
        \mathcal C^{A^0}_{c,r}M
         & =|G|^{-|X|}\sum_{q:X\to G}\mathcal C^{B^q}_{c,r}M.
        \label{eq:peps_common_cut_average_expansion}
    \end{align}
    This includes the matched matrices
    (\ref{eq:peps_common_matched_action}). On the two-by-two torus the
    normalization is $|G|^{-4}$ and there are four independent group labels.
\end{theorem}
\begin{proof}\leanok
    Distribute the finite sums independently at every site and exchange
    them with the bond-end sum. Site scalars factor out as their product.
    The averaging formula gives $|G|^{-1}$ at each site and the stated sum
    of matrix entries. Throughout, the same function $M$ multiplies each
    bond-end configuration.
\end{proof}

\begin{definition}[Independent four-leg changes of basis]%
    \label{def:peps_four_cut_common_leg_gauge}
    \lean{TNLean.PEPS.torusLegGaugeMatrix,
        TNLean.PEPS.torusLegGaugeMatrix_apply,
        TNLean.PEPS.torusLegDress,
        TNLean.PEPS.torusLegDress_single}
    \leanok
    For four matrices $T,R,B,L$ on $\C^V$, define
    $K=T\otimes R^{\mathsf T}\otimes B^{\mathsf T}\otimes L$ and,
    for a row tensor $a:V^4\to\C$, define its dressed row to be $aK$.
    In coordinates,
    \begin{align}
        K_{c,c'} & =T_{c_1,c'_1}R_{c'_2,c_2}B_{c'_3,c_3}L_{c_4,c'_4},
        \nonumber                                                     \\
        (aK)(c') & =\sum_{c\in V^4}a(c)K_{c,c'}.
        \nonumber
    \end{align}
    For the fixed row $a(c)=\delta_{s,c}$, the dressed row is $K_{s,c'}$.
\end{definition}
\begin{proof}\leanok
    The entries follow from the tensor-product matrix and transpose
    formulas. The fixed row removes every summand except $c=s$.
\end{proof}

\begin{theorem}[Moving independent leg operators across a bond]%
    \label{thm:peps_four_cut_common_leg_gauge}
    \lean{TNLean.PEPS.torusBondNetwork_legGauge}
    \leanok
    \uses{def:peps_four_cut_common_leg_gauge,def:peps_torus_bond_network}
    For arbitrary site rows $a_v$ and matrix families
    $O^{\mathrm h},O^{\mathrm v},T,R,B,L$, put
    $K_v=T_v\otimes R_v^{\mathsf T}\otimes B_v^{\mathsf T}\otimes L_v$.
    Then the scalar bond contractions satisfy
    \begin{align}
        \mathcal N_a((L_{v+e}O^{\mathrm h}_vR_v)_v,
        (T_vO^{\mathrm v}_vB_{v+n})_v)
         & =\mathcal N_{(a_vK_v)_v}(O^{\mathrm h},O^{\mathrm v}).
        \label{eq:peps_common_leg_gauge}
    \end{align}
\end{theorem}
\begin{proof}\leanok
    Expand each three-matrix bond product and each dressed site row.
    The bijection $\beta\mapsto(s_v(\beta))_v$ from
    (\ref{eq:peps_common_boundary_restriction}) regroups the intermediate
    endpoint indices by site. Reindex the factors attached to $L_{v+e}$
    and $B_{v+n}$ by translation. The scalar summands on both sides then
    coincide, with the same original bond matrices remaining between them.
\end{proof}

\begin{theorem}[Matched vertex gauge invariance]%
    \label{thm:peps_four_cut_common_vertex_gauge}
    \lean{TNLean.PEPS.vecMul_torusMatchedLegMatrix_of_comp_eq,
        TNLean.PEPS.torusBondNetwork_matchedVertexGauge}
    \leanok
    \uses{def:peps_four_cut_common_matched_action,def:peps_torus_bond_network}
    Suppose $\mathcal P(A_v)\rho_v(k)=\mathcal P(A_v)$ for every $v,k$.
    Each physical row of $A_v$ therefore satisfies $A_v^{s}\rho_v(k)=A_v^{s}$.
    For every $q:X\to G$, put
    \begin{align}
        O^{\mathrm h,q}_v
         & =U^{\mathrm h}_v(q_{v+e})O^{\mathrm h}_vU^{\mathrm h}_v(q_v^{-1}),
        \label{eq:peps_common_gauged_horizontal}                              \\
        O^{\mathrm v,q}_v
         & =U^{\mathrm v}_v(q_v)O^{\mathrm v}_vU^{\mathrm v}_v(q_{v+n}^{-1}).
        \label{eq:peps_common_gauged_vertical}
    \end{align}
    Then $\mathcal N_A(O^{\mathrm h,q},O^{\mathrm v,q})
        =\mathcal N_A(O^{\mathrm h},O^{\mathrm v})$.
\end{theorem}
\begin{proof}\leanok
    \uses{thm:peps_four_cut_common_leg_gauge}
    Test the site-map invariance on each virtual basis vector to obtain
    the row identity. In (\ref{eq:peps_common_leg_gauge}), choose
    $T_v=U^{\mathrm v}_v(q_v)$,
    $R_v=U^{\mathrm h}_v(q_v^{-1})$,
    $B_v=U^{\mathrm v}_{v-n}(q_v^{-1})$, and
    $L_v=U^{\mathrm h}_{v-e}(q_v)$.
    These four factors form $\rho_v(q_v)$, which each physical row absorbs.
\end{proof}

\begin{theorem}[Matched averaging networks as coherent bond products]%
    \label{thm:peps_four_cut_common_network_expansion}
    \lean{TNLean.PEPS.torusBondNetwork_representationAveragingSite,
        TNLean.PEPS.torusBondRegrouping_insertedAveragingSite_coherent,
        TNLean.PEPS.torusGaugeBondMatrix}
    \leanok
    \uses{def:peps_four_cut_common_matched_action,
        def:peps_torus_physical_bond_regrouping}
    Let $G$ be finite and use the bond matrices
    (\ref{eq:peps_common_gauged_horizontal}) and
    (\ref{eq:peps_common_gauged_vertical}). Then
    \begin{align}
        \mathcal N_{A^0}(O^{\mathrm h},O^{\mathrm v})(\sigma)
         & =|G|^{-|X|}\sum_{q:X\to G}\prod_{v\in X}
        O^{\mathrm h,q}_v(l_{v+e},r_v)
        O^{\mathrm v,q}_v(t_v,b_{v+n}),
        \label{eq:peps_common_network_average}                                   \\
        (\mathcal R\mathcal N_{A^0}(O^{\mathrm h},O^{\mathrm v}))(\beta)
         & =|G|^{-|X|}\sum_{q:X\to G}\prod_{f\in X\times\{\mathrm h,\mathrm v\}}
        O^q_f(\beta_f^1,\beta_f^2).
        \label{eq:peps_common_regrouped_average}
    \end{align}
    Here $\sigma_v=(t_v,r_v,b_v,l_v)$, $\mathcal R$ is physical bond
    regrouping, and $O^q_{(v,d)}=O^{d,q}_v$ defines the single bond-indexed
    family. No semi-regularity or unitarity is required.
\end{theorem}
\begin{proof}\leanok
    \uses{thm:peps_four_cut_common_cut_expansion,
        thm:peps_four_cut_common_leg_gauge,thm:peps_torus_projector_expansion,
        def:peps_torus_physical_bond_regrouping,
        lem:peps_torus_physical_bond_regrouping_overlap}
    Expand each local average as a sum of its matched matrices. Finite
    multilinearity produces $|G|^{-|X|}\sum_q$.
    For each $q$, apply (\ref{eq:peps_common_leg_gauge}) to the fixed
    rows $\delta_{\sigma_v}$; the remaining identity-site contraction
    is the product of its bond entries. This gives
    (\ref{eq:peps_common_network_average}). Substitute the inverse
    physical regrouping formulas and combine the horizontal and vertical
    products to obtain (\ref{eq:peps_common_regrouped_average}).
\end{proof}
